\subsection{Graphing General Logarithms} 
Graphing a general logarithm is not \emph{difficult}, but it combines two new strategies.
\begin{itemize}
\item We want to look at the \makeblank[1in] and \makeblank[1in] of the \emph{base function} instead of the given function.
\item Because we are graphing an \makeblank[1in] we want to choose values of the \makeblank[1in].
\end{itemize}
\begin{example}
Sketch the graph of $2\log_{1/3}(4-x)-3$.
\begin{cpic}{}{16}
	\draw (-5,-2) -- (5,-2);
	\foreach \x in {-3,0,3}
		\draw (\x,-1) -- (\x,-7);
	\foreach \x/\lbl in {-4/$x$,-1.5/input,1.5/output,4/$y$}
		\regtext{\x}{-2}{\lbl};
\end{cpic}
\end{example}
We can see the key features of this graph, which will help us graph more easily in the future.
\begin{itemize}
\item The vertical asymptote is where \makeblank equals \makeblanksmall.
\item If the \makeblank in the \makeblank[1in] is negative (as in this example) the graph is to the \makeblank[1in] of the asymptote.
\item To find nice key points to draw our graph, we can start with the \makeblank[1in] $-1$, $0$, and $1$, which correspond to \makeblank[1in] of \makeblanksmall, \makeblanksmall, and \makeblanksmall. This will usually give us two nice points and one fraction.
\end{itemize}